% PROJECT DESCRIPTION -- 15-page limit.
\section{Intellectual merit}
\label{sec:IM}

Ice-sheet science increasingly depends on integrating nonlinear numerical
models, satellite, airborne, radar, and in-situ observations, ensemble data
assimilation (DA), and scalable computing. DA provides a principled means to
update model state and uncertain parameters while representing forecast uncertainty
\citep{evensen2003,gilletchaulet2020,choi2025}. In practice, however, an
end-to-end DA experiment remains a systems-integration project. A scientist
must install a model and its compiled dependencies, translate fields between
model and DA representations, stage observations, configure MPI, submit and
monitor scheduler jobs, recover outputs, and preserve enough metadata to
reproduce the calculation. These tasks are repeated for every model and
computing site. They divert scientific effort, discourage cross-model studies,
slow collaboration, and exclude groups that lack dedicated research-software
engineering or HPC support. The underlying problem is broader than any one
model or algorithm: community applications repeatedly rebuild authentication,
configuration, storage, visualization, provenance, and remote-execution
infrastructure instead of sharing sustainable services.

The project team has established exploratory components that demonstrate
feasibility but are not yet community cyberinfrastructure. \textbf{ICESEE} is
a Python/MPI ensemble DA research framework with demonstrated ISSM and Icepack
couplings \citep{kyanjo2026icesee}, executed through its own Local, Remote,
and Cloud implementations. \textbf{CryoLauncher} configures and runs ISSM and
Icepack through Remote (Direct SSH or Connector/Relay to institutional
Slurm/PACE resources) and Cloud (AWS Batch on Fargate or EC2 On-Demand)
execution paths, with no local execution mode, and is being decomposed into
reusable frontend panels, execution backends, remote drivers, and
cloud-provider interfaces. A shared
gateway now provides salted-password, GitHub, and ORCID authentication;
schema-versioned configurations and workspaces; experiment records and event
timelines; and a role-based Control Center for users, experiments,
authentication, diagnostics, and audit events. \textbf{LIVIST} (Living Ice
Sheet Temperature) explores Antarctic englacial-temperature products inferred
from radar and constrained by boreholes. The new \textbf{Frozen Legacies}
application registers historical Antarctic radar collections through dataset
manifests and adapters, builds catalog and GeoJSON products, and exposes
flightlines, observation points, derived LYRA products, and the ASTRA, ARIES,
TERRA, and URSA interpretation tools. CryoLauncher and ICESEE implement their
own submission, execution, and result-handling logic rather than a shared
pipeline, but reuse the same Connector/Relay connectivity and a
bring-your-own-AWS onboarding path -- connected through connection-scoped
CloudFormation, so independently created connections do not depend on one
globally fixed execution role -- instead of duplicating separate cloud
stacks. On that shared onboarding path, CryoLauncher's ISSM workflow has
completed end-to-end execution on both AWS Batch Fargate and EC2 On-Demand,
including MPI-parallel solver execution, postprocessing, result retrieval,
and visualization, and has reached the Georgia Tech institutional MATLAB
license through the same Connector/Relay infrastructure used for Remote
access; ICESEE's own Cloud path has separately executed its Lorenz-96
example end to end on Fargate at a single process (NP=1). Together these
components establish technical feasibility for a shared data-to-model
environment \citep{cryostack2026}; demonstrated end-to-end execution does not
establish a scientifically qualified or secure cloud lifecycle, and the
components do not yet provide a stable community plug-in contract,
production-safe HPC bridge, qualified cloud lifecycle, complete provenance,
or distributed operations model.

Operational evidence from the current deployment complements this technical
feasibility. The four applications run behind one gateway that shares
identity, roles, workspaces, and the Control Center while each keeps its own
scientific interface and data or execution path; a dependency-aware registry
builds the Frozen Legacies, LIVIST, and documentation components, and health
checks cover the hosted applications. The October 2026 release added
login-protected application routes, gateway-only access to the notebook
servers, per-user grants for internal and relay requests, an allow-listed
notebook environment that excludes gateway credentials, and a backup procedure
verified by restoring production data into an isolated copy; its working tree
passes 3,649 automated tests. This evidence describes a deployed and tested
working tree, not yet a committed or archived source release, and it does not
remove the limitations below.

CryoStack is designed as general scientific cyberinfrastructure rather than a
wrapper around these initial applications. Scientific tools retain their
domain-specific interfaces and capabilities while reusing identity,
configuration, experiment management, execution, storage, visualization,
deployment, and provenance services. This award commits to PISM as the new
production model adapter and to a second Frozen Legacies dataset adapter as a
test of data-side extensibility. It establishes the contracts through which
future models, data collections, preprocessing utilities, and analysis tools
can enter without redesigning the platform. Other models such as Elmer/Ice and
SICOPOLIS remain a forward pathway, not award deliverables.

These components constitute a credible prototype, not yet sustained
cyberinfrastructure. The modular migration is incomplete; application and
dataset contracts are not versioned public SDKs; tests do not form a common
scientific acceptance gate; connector sessions are bound to their owner but
not to resource policy or a durable session store; broad remote commands and
an optional password-bearing bootstrap remain; notebook kernels still run under
the gateway's operating-system account, and the per-user kernel sandbox is
tested but not enabled in production; production state uses SQLite, and the
health checks are not an external availability monitor; the demonstrated CryoLauncher AWS workflow is not yet
CryoBench-qualified for production use and EC2 Spot, GPU, and multi-node
provisioning remain scientifically unvalidated; metadata and licenses
disagree across packages; and adoption is not yet measured against public
targets.
\textbf{The central premise of the proposed project is that} scientific
applications should be able to evolve independently while reusing secure,
portable, and measurable services for identity, experiment management,
execution, provenance, storage, and visualization. Achieving this requires
research and engineering beyond routine operation of the prototype:
generalizing interfaces, replacing unsafe remote-control patterns, qualifying
science across heterogeneous resources, evaluating users, and distributing
maintenance beyond one developer. \textbf{The project tests the hypothesis
that} models and observational datasets with different interfaces, formats,
and scientific uses can share one execution, provenance, and access
infrastructure without losing domain-specific scientific meaning, provided
that interoperability is defined by versioned, scientifically validated
contracts rather than by convention, and that doing so reduces the effort
required to reproduce, extend, and sustain cryosphere workflows. Three
integrated objectives test this hypothesis: O1 establishes scientifically
meaningful interoperability, O2 qualified and reproducible execution, and O3
reproducible use, evaluation, and sustained access:

\mybox{
\vspace{-3pt}
\begin{enumerate}[topsep=0pt,leftmargin=*]
  \item[\textbf{O1}] \textbf{Develop and validate scientifically
  interoperable and extensible infrastructure for cryosphere models and
  observational datasets.} Complete the modular platform migration; define versioned
  contracts and templates; deliver PISM and a second historical-radar dataset
  adapter; and build \textbf{CryoBench} qualification gates.
  \item[\textbf{O2}] \textbf{Develop and evaluate reproducible, secure, and
  resource-aware scientific execution across heterogeneous computing
  environments.} Develop a typed FastAPI control plane and least-privilege
  connector/relay; replace unrestricted remote execution with allow-listed
  operations; implement least-privilege credential handling, per-user kernel
  isolation, policy-controlled computing allocations, and audit trails; package reproducible
  environments; and qualify consistent workflows on Georgia Tech PACE and a
  user- or project-managed AWS deployment.
  \item[\textbf{O3}] \textbf{Establish and evaluate reproducible scientific
  experiment services and sustainable community-access mechanisms for
  cryosphere research.} Extend the existing identity, experiment, Control
  Center, and deployment foundations into durable provenance and operations
  services; connect Frozen Legacies or LIVIST products to modeling/DA; evaluate
  all four applications through representative scientific tasks; and
  establish contribution, release, access, and governance processes.
\end{enumerate}
\vspace{-3pt}
}

The requested three-year Elements award is the appropriate scale: it will
harden and generalize an early research prototype while focusing the new
scientific integration on one high-value community model, PISM, and one new
historical-radar collection. Every future application or dataset adapter will
reuse the same identity, validation, execution, provenance, deployment, and
operations services, and every qualified backend will serve all executable
applications. This gives CryoStack high growth potential: the marginal cost of
future integrations falls while the scientific and preservation value of the
shared foundation grows. It will not create a new ice-sheet model,
operate a general-purpose commercial cloud, or promise permanent compute to
all users.

\section{State of knowledge}
\label{sec:SOK}

\subsection{Fragmented scientific software and the missing community layer}

ISSM, Icepack, and PISM embody different numerical methods, languages,
dependency stacks, restart conventions, and execution assumptions
\citep{larour2012,shapero2021,pism2026}. This diversity is scientifically
valuable, but each surrounding workflow separately implements installation,
configuration, state translation, observation handling, job submission, and
result packaging. Ensemble DA adds deterministic member initialization,
state packing and unpacking, observation operators, synchronized failure
handling, parallel I/O, and diagnostics. Current tools do not provide a common
path by which an independently developed model can declare these capabilities,
pass scientific and operational qualification, and become portable across
resources. CryoStack's new backend interfaces and frontend panels establish
the code boundaries for such a path, but still wrap portions of the original
monolithic implementations.

The same gap separates observations from simulation. LIVIST serves derived
temperature products, while Frozen Legacies now demonstrates manifest-driven
registration, ingestion adapters, catalog/GeoJSON generation, and preservation
of historical radar interpretation tools. Its review workflow records the stage
and version of each submitted interpretation, and automatically generated URSA
picks carry source, parameter, and per-frame quality provenance while remaining
marked unvalidated; LIVIST records the upstream identity of each input and
compares generated products by checksum. These data applications need
provenance and review semantics that a model-execution interface does not
express, which is why dataset and model contracts must be distinct. Only the
initial LYRA adapter is registered, however, and desktop processing is not yet a browser-executable,
provenance-captured workflow. Without versioned dataset and
observation-bundle contracts, units, coordinates, uncertainties, quality
controls, transformations, source identifiers, and licenses can be lost
during handoff. A community layer should preserve specialized interfaces
while making exchanged artifacts explicit and reproducible.

\subsection{Heterogeneous computing and science-gateway execution}

Scientific workflows must run where allocations, data, licensed software,
and suitable hardware exist. Institutional clusters often use Slurm and
restrict inbound traffic; cloud systems require cost, identity, quota, and
lifecycle controls. Containers and Spack improve environments but do not supply job
authorization, staging, monitoring, cancellation, recovery, or a portable
record of execution. The unresolved challenge is a narrow bridge between
a public control plane and protected, user-authorized resources. Storing
long-lived private keys or exposing arbitrary shell commands is unacceptable;
requiring users to rebuild scheduler and provenance logic at each site defeats
the gateway's purpose.

\subsection{Reproducibility, qualification, and proof of concept}

Source code alone does not reproduce an experiment. The record must bind code
and adapter versions, configuration, inputs and transformations, environment
digests, resource request, random seeds, scheduler identity, logs, outputs,
checksums, and diagnostics. Nor does a zero exit code establish scientific
validity: conservation, restart consistency, assimilation error and spread,
and model-specific tolerances require explicit tests.

ICESEE, CryoLauncher, LIVIST, Frozen Legacies, authentication and experiment
services, the Control Center, packaging recipes, and the deployment registry
provide a proof of concept for a shared environment. They also expose the
remaining gap. Interfaces are not yet stable external SDKs; shared platform
tests do not yet gate scientific or cross-backend execution equivalence; the
connector command channel is too broad for public use; notebook kernels are not
yet isolated from the gateway's operating-system account;
credential, relay, and recovery semantics have not received external review;
PISM is not integrated; the CryoLauncher AWS workflow executes end to end
for ISSM and Icepack but is not yet CryoBench-qualified, and ICESEE's own
Cloud path remains limited to a single-process Lorenz-96 demonstration;
provenance is not yet a stable cross-application schema; and adoption is not
measured against public targets.
The proposed award is therefore not funding an already completed portal. It
funds the distinct work required to convert promising components into trusted,
extensible community infrastructure.

\subsection{Scientific opportunity and community impact}

The central science outcome is a reproducible path from observations to
multi-model estimates of ice-sheet state and uncertain parameters. PISM is an
open-source, parallel, thermodynamically coupled ice-sheet model that has been
used in cryosphere research for roughly two decades \citep{pism2026}. Its
addition gives CryoStack a complementary model architecture and a direct path
to a mature user community. The project will support three motivating use
cases:

\begin{itemize}[leftmargin=*]
  \item \textbf{PISM state and parameter estimation.} Assimilate synthetic
  and observational surface velocity and elevation products to estimate geometry
  and selected basal or rheological parameters in a regional Greenland test
  domain. The use case will include an observing-system simulation experiment
  with known truth and a documented real-data demonstration.
  \item \textbf{Cross-model reproducibility.} Run a common observation bundle,
  ensemble specification, and diagnostic definition through two or more
  supported adapters. Differences can then be attributed to model structure
  or configuration rather than undocumented staging and post-processing.
  \item \textbf{Data-to-model workflow.} Define a versioned observation-bundle
  schema and demonstrate the path from a Frozen Legacies radar-derived product
  or LIVIST temperature product to model initialization, validation, or an
  assimilation experiment.
  The schema will record variable definitions, coordinates, units,
  uncertainty, source DOI/URL, spatial/temporal coverage, and transformations.
\end{itemize}

A parallel preservation use case will register one additional historical
radar collection through the Frozen Legacies manifest/adapter path, record
source and transformation provenance, and publish machine-readable catalog,
geospatial, quality-control, and downloadable derived products. This tests
whether the same platform can grow through data contributions as well as model
contributions.

The work benefits users beyond the initial PISM target because the same
adapter, validation, provenance, and execution contracts apply to glacier,
climate, and geophysical models. Tutorials will be usable on a laptop with an
idealized model before a user requests HPC access, enabling education and
participation without making training the primary project activity.

\subsection{Innovation required beyond the prototypes}

The innovation is not a new web skin around existing codes. CryoStack will
make scientific interoperability an executable and testable contract across
applications, models, datasets, and compute resources. A model adapter will
expose: model identity and version; configuration schema;
initialization and forecast operations; state-vector packing/unpacking;
observation operators; restart/checkpoint behavior; resource estimates; and
output/provenance declarations. Adapters will be isolated from execution
backends, so the PISM integration will not contain PACE- or AWS-specific job
logic. Conversely, execution backends will submit a typed experiment without
knowing the internal details of PISM or ISSM.

This separation makes the experiment specification portable across resources.
The current prototype already demonstrates this informally: the same ISSM or
Icepack configuration submitted through CryoLauncher runs on institutional
Slurm via Remote execution or on AWS Batch via Cloud execution without a
backend-specific rewrite of the scientific application, and ICESEE
separately demonstrates the same property across its own Local, Remote, and
Cloud implementations. What does not yet
exist is a \emph{formal, versioned} contract that makes this property hold
for an independently developed model such as PISM by construction rather than
by convention: a user can move a validated workflow from a workstation to an
institutional cluster or a cloud account supported by the user's own
allocation or credits without rewriting the scientific application. Conversely,
a newly integrated model can reuse authentication, workspaces, experiment
history, remote execution, result delivery, and visualization instead of
developing parallel infrastructure. The reduction of duplicated delivery
software is central to CryoStack's sustainability model.

CryoBench will turn this contract into an acceptance gate. Unit and schema
tests will check interfaces; manufactured/analytic and restart tests will
check model behavior; seeded DA tests will compare error and spread
diagnostics against tolerances; scaling tests will record cost; software bills
of materials and vulnerability scans will check deployability; and a complete
provenance manifest will be required before release. CryoBench will also
qualify \textbf{cross-backend scientific equivalence}: the question is not
whether a workflow merely executes on Local, PACE, and AWS, but whether the
same experiment remains scientifically equivalent across them. Equivalence
will be defined by application-appropriate, versioned tolerance envelopes
rather than bitwise reproducibility --- field norms, integrated quantities,
conserved or diagnostic quantities, ensemble statistics, and DA-specific
diagnostics such as error and spread, depending on the model and experiment.
This is deliberately distinct from performance portability: CryoBench will
report, but never require, matching runtime or cost across systems, only that
a qualified backend produces a scientifically acceptable answer. Dataset
adapters will declare source identity, schema, coordinates, units, uncertainty
and quality fields, transformations, licenses, products, and validation
fixtures. This is transformative for the target community because a model
becomes a reliable, launchable
service only after its scientific and operational assumptions are explicit.

\section{Proposed research}
\label{sec:research}

\subsection{System Architecture}

\begin{figure}[ht]
\centering
\setlength{\fboxsep}{7pt}
\begin{tabular}{c}
\fbox{\parbox{0.88\textwidth}{\centering
\textbf{Users and contributors} (existing) --- browser, notebooks, documented APIs}}\\[3pt]
$\downarrow$ HTTPS\\[-1pt]
\fbox{\parbox{0.88\textwidth}{\centering
\textbf{Gateway and operations plane} (existing) --- federated identity, RBAC,
workspaces, experiments/events, Control Center, deployment registry, audit}}\\[3pt]
$\downarrow$ versioned application and data contracts (\textit{proposed, O1})\\[-1pt]
\fbox{\parbox{0.88\textwidth}{\centering
\textbf{Scientific applications} (existing) --- CryoLauncher models $\mid$ ICESEE DA $\mid$
LIVIST temperature products $\mid$ Frozen Legacies radar data/tools}}\\[3pt]
$\downarrow$ typed job specification + immutable environment reference (\textit{proposed, O2})\\[-1pt]
\fbox{\parbox{0.88\textwidth}{\centering
\textbf{Connector and relay} (existing: outbound authenticated session, and a
demonstrated Cloud-to-institutional-service bridge for CryoLauncher/ISSM;
\textit{proposed:} signed job envelopes, allow-listed
stage/submit/query/cancel/recover operations)}}\\[3pt]
$\downarrow$ backend adapter; bulk data uses an approved transfer path\\[-1pt]
\fbox{\parbox{0.88\textwidth}{\centering
\textbf{Execution resources} (existing options; \textit{proposed:} qualified PACE and
bring-your-own-AWS profiles) --- local (ICESEE) $\mid$ institutional
Slurm/PACE via Direct SSH or Connector/Relay (Remote) $\mid$ AWS Batch
Fargate or EC2 On-Demand (Cloud)}}\\[3pt]
$\downarrow$\\[-1pt]
\fbox{\parbox{0.88\textwidth}{\centering
\textbf{Models, data, and environments} (existing: ISSM, Icepack; \textit{proposed:} PISM) ---
historical radar and temperature products; NetCDF/HDF5/Zarr; Spack/OCI/Apptainer}}\\[2pt]
$\uparrow$ signed status, logs, outputs, and provenance manifest (\textit{proposed contract, O3})
\end{tabular}
\caption{CryoStack's existing foundation and the CSSI advances proposed on
top of it. The underlying CryoStack layers shown are present in the current
prototype; \textit{italicized, proposed} annotations identify capabilities
that will be generalized, hardened, added, or scientifically qualified
through this award. The relay will carry typed control messages and state,
not arbitrary shell access or bulk scientific data, and credentials will
remain at the narrowest user- or site-controlled layer.}
\label{fig:architecture}
\end{figure}

Relative to this existing foundation, the award funds three concentrated
advances. \textbf{O1} replaces ad hoc model/data integration with versioned
extension contracts, independently tested by PISM and a second Frozen
Legacies collection and gated by CryoBench's scientific and cross-backend
equivalence tests. \textbf{O2} replaces the connector's broad shell-command
channel with the typed, allow-listed control plane shown above and qualifies
PACE and bring-your-own-AWS execution, including representative parallel
workloads where the science requires them. \textbf{O3} replaces single-node,
SQLite-backed state with a durable, versioned experiment/result/provenance
contract, evaluated with real users and sustained through public governance.

The proposed FastAPI service will expose versioned endpoints whose
Pydantic/JSON schemas define experiment creation, validation, authorization,
submission, cancellation, artifact registration, and state transitions. Its
OpenAPI contract will make the control boundary testable; scientific adapters
will remain independent of the web framework. Each experiment receives a unique identifier and
an immutable configuration snapshot. A job specification references an
adapter version, observation-bundle checksum, container digest or Spack lock
file, resource request, random seed(s), and requested diagnostics. Status
transitions, scheduler job ID, backend, timestamps, exit code, logs, and output
locations are recorded. On completion, a JSON manifest binds these fields to
code commits, image digests, dependency inventory, and input/output checksums.

\subsection{Qualified application, model, and data extension kits (\textbf{O1})}

\textbf{Task 1.1 -- Extension specifications and modular migration.} The O1
research question is whether versioned contracts can capture the state,
configuration, provenance, and validation needs of models and datasets that
differ in interfaces, formats, computational requirements, and scientific use,
without erasing their application-specific meaning. Model contracts describe
state, forecast, restart, and observation-operator behavior; dataset contracts
describe source identity, coordinates, uncertainty, review status, and product
lineage. The contribution is a tested method for extending scientific
cyberinfrastructure, evaluated by PISM (Task 1.2) and a second historical-radar
collection (Task 1.4) rather than by the components that shaped it. During
months 1--9, the team will finish separating the current CryoLauncher panels,
state objects, execution managers, remote drivers, cloud drivers, and
container services from legacy gateway callbacks. Versioned application,
model, dataset, and observation-bundle contracts will use typed Python
interfaces and JSON Schema, semantic versioning, deprecation rules, reference
fixtures, and template repositories. A compatibility matrix will state which
application, model/data adapter, environment, and backend combinations are
supported. The dependency-aware deployment registry will become the common
application packaging and health contract.

\textbf{Task 1.2 -- PISM integration as an independent test of the extension
architecture.} CryoLauncher, ISSM, Icepack, and ICESEE helped motivate
Task 1.1's emerging adapter interfaces, so their fit is not a strong test of
generality. PISM did not shape those interfaces: it is an independently
developed model with its own state representation, I/O, and restart
conventions. Its integration therefore tests whether Task 1.1's versioned
contracts generalize to a model the team did not design around, not merely
whether one more model can be added. During months 4--18, the team will wrap
PISM restart and NetCDF I/O, map selected state/parameter fields, implement
observation operators, and support deterministic ensemble-member execution
using only the adapter's public interface. The team will classify every
integration change by code review: changes confined to the adapter and its
declarative/configuration layer count toward a target of at least 90\%, and
any remaining change to shared CryoStack code must enter through a
documented extension interface rather than a model-specific branch or
special case, as a direct, auditable measure of whether the extension kit
reduces integration burden. The proof of concept
will reproduce a documented PISM regional or idealized case and a
synthetic-truth assimilation. A production release requires conservation/restart
checks, expected error reduction without ensemble collapse, reproducibility
within stated tolerances, and successful, scientifically equivalent runs on
two qualified execution backends. The team will seek design review from PISM
developers or experienced users; no changes to PISM itself are required for
core delivery.

\textbf{Task 1.3 -- CryoBench.} During months 6--24, the team will implement a
tiered validation pipeline: fast checks on every pull request; model-specific
CPU tests on merge; scheduled MPI/container tests; and release-candidate tests
on supported backends. The release-candidate tier is where cross-backend
scientific equivalence is enforced: a candidate adapter/backend combination
must reproduce its reference tolerances on every qualified backend it claims
to support, not merely execute on it. Results, reference data, tolerances,
runtime, and software versions will be published. A failed required tier
blocks promotion to the CryoStack production registry. By month 30, every
supported adapter will have an identified maintainer, test owner,
compatibility statement, and documented scientific limitations.
CryoBench criteria are application-appropriate. Computational adapters are
judged on numerical consistency, scientific equivalence across backends,
configuration and environment reproducibility, and recorded performance and
resource use; data adapters are judged on identifiers, coordinates and units,
checksums, provenance completeness, and review status rather than on
simulation benchmarks. Computing cost is recorded as a run attribute, not used
as a qualification criterion.

\textbf{Task 1.4 -- Frozen Legacies extension and data qualification: an
independent test of dataset extensibility.} Where Task 1.2 tests whether the
model-adapter contract generalizes to an independently developed model, this
task tests whether the parallel dataset-adapter contract generalizes beyond
the single LYRA adapter it was designed around. During
months 7--24, the project will turn the existing manifest, adapter, catalog,
GeoJSON, map, and documentation implementation into a documented dataset SDK.
One additional historical radar collection will be ingested. Qualification
will validate identifiers, coordinates, units, coverage, checksums, licenses,
source and transformation provenance, quality fields, downloadable products,
and geospatial continuity. At least one ASTRA, ARIES, TERRA, or URSA processing
path will be wrapped as a reproducible noninteractive task with captured
parameters and QC outputs, while retaining interactive tools for expert review.

\subsection{Reproducible, secure, and resource-aware execution (\textbf{O2})}

The O2 research question is whether a declared scientific experiment can keep
its scientific meaning and provenance when security, isolation, and resource
policies are enforced by systems as different as an institutional cluster and a
cloud account. The proposed work is not to build an HPC connector; a connector already runs
on a user workstation, login/edge host, or institution-approved service near
the compute resource, and already establishes an outbound authenticated
session to the relay, avoiding any requirement for inbound access through an
institutional firewall. The same connector/relay infrastructure already also
lets a Cloud workload reach a private institutional service when required --
the demonstrated case is CryoLauncher's ISSM workflow on AWS reaching the
Georgia Tech institutional MATLAB license -- which validates one
institutional Cloud/Connector configuration rather than a general
private-service bridging capability. The proposed work is to transform that existing
proof-of-concept connector into a constrained, auditable control plane
suitable for bridging a public gateway to protected institutional resources:
the relay will queue signed, typed job envelopes and return state events
instead of exposing an interactive shell, and large inputs and outputs will
use a configurable approved transfer path rather than passing through the
control relay.

\textbf{Task 2.1 -- FastAPI control plane, threat model, and connector hardening.} The current desktop
connector proves that a browser workflow can reach user-authorized HPC
resources, but its general shell-command channel is too broad for a production
service. In months 1--9, the project will document assets, trust boundaries,
actors, misuse cases, and incident procedures. The connector protocol will be
replaced with allow-listed, schema-validated operations (stage declared files,
submit a generated batch script, query/cancel an owned job, fetch declared
outputs). Commands will not accept arbitrary shell strings. Paths will be
canonicalized into per-experiment directories; resource limits and input sizes
will be enforced; logs will redact secrets; and dependencies/images will be
pinned and scanned.

Private keys and passwords will never be stored in the public gateway
database. For campus resources, credentials will remain in the user's local
connector or an institution-approved credential service. Connector sessions
will be bound to authenticated users, registered resource profiles, explicit
capabilities, expiry, and revocation; command envelopes will include ownership,
nonce, and replay checks.
The gateway already enforces login-protected routes, gateway-only notebook
access, per-user grants, and an allow-listed notebook environment, covered by
automated regression tests and a rehearsed restore. The project will extend
these tests to session rotation, CSRF protection, rate limits, patching,
vulnerability scanning, and incident response, and will isolate notebook
kernels in two layers: a per-user sandbox that exposes only the user's own
workspace (implemented and tested, not yet enabled in production) and a
separate service account for the notebook servers. An acceptance test will
verify that a kernel cannot reach gateway credentials, other users'
workspaces, or operator keys. An independent Georgia Tech security review will
be requested by month 18.

\textbf{Task 2.2 -- Reproducible environments.} The team will consolidate the
existing ICESEE-Containers and ICESEE-Spack recipes into supported profiles.
Each release will publish source provenance, lock files, OCI/Apptainer image
digests, software bills of materials, license notices, and smoke-test results.
The Year-1 release will reconcile the present BSD/MIT metadata discrepancy and
document third-party license constraints, including MATLAB-dependent ISSM
paths. Open PISM/Icepack paths will provide the default no-proprietary-license
tutorial experience.

\textbf{Task 2.3 -- Backend qualification.} Qualification deliberately spans
two different infrastructure regimes: PACE, an institutionally managed HPC
environment with scheduler, identity, policy, and protected-resource
constraints; and AWS, a researcher- or project-controlled cloud environment
with its own IAM, Batch, object storage, quota, cost, and lifecycle concerns.
Both will be qualified through the same scientific experiment and CryoBench
acceptance framework, without requiring identical implementations: the
scientific question is whether a qualified workflow retains valid,
reproducible scientific meaning across these heterogeneous environments, not
whether the environments themselves behave alike. PACE remains the primary
development and science-validation resource. AWS is used for portability and
elastic demonstration, with budget and quota controls. CryoStack does not
require AWS: Slurm/PACE execution remains fully supported, and AWS bills cloud
computing to the account that runs it -- by default the researcher's own
account, and the project account only for the bounded qualification and access
activities budgeted here. Building on the
existing bring-your-own-AWS path -- which already connects a researcher's own
account through connection-scoped CloudFormation onboarding, so independently
created connections do not depend on one globally fixed execution role, and
on which CryoLauncher's ISSM and Icepack workflows have already completed
end-to-end execution on both AWS Batch Fargate and EC2 On-Demand, including
MPI-parallel ISSM solver execution, postprocessing, result retrieval, and
visualization, while ICESEE's own Cloud path has separately completed a
single-process (NP=1) Lorenz-96 run on Fargate -- the project will qualify
this path for sustained scientific use: per-user object isolation in place of
a single account-wide bucket, allow-listed Batch job definitions in place of
a caller-supplied one, budgets, quotas, lifecycle cleanup, failure recovery,
cost attribution, and least-privilege credential handling, tested against a
real AWS account rather than fakes alone. These verified results are useful
preliminary evidence precisely because they expose the boundary between
architectural capability and scientific qualification: CryoLauncher's ISSM
and Icepack workflows and ICESEE's Lorenz-96 workflow already execute end to
end on AWS, but qualification, not new execution architecture, is what makes
a configuration trustworthy for production use. CryoLauncher's ISSM workflow
already demonstrates single-node MPI-parallel execution on Cloud; where a
scientific use case requires it, this capability will be scientifically
qualified through CryoBench and extended to multi-node execution and to
ICESEE's own multi-process ensemble workflows on Cloud. The project does not
commit to a general-purpose cloud-MPI platform independent of the scientific
use cases it serves. A backend is
``qualified'' only if it passes staging, submission,
monitoring/cancellation, output recovery, provenance, failure-recovery, and
the cross-backend scientific-equivalence tests defined in CryoBench
(Task 1.3). The project will not claim identical performance across
systems; resource profiles and measured cost/time will guide users.

\textbf{Task 2.4 -- Policy-controlled computing allocations.} CryoStack will
support two independent AWS execution models -- a researcher's own account,
billed by AWS to that account, and CryoStack-managed computing funded by the
project -- while institutional HPC execution remains independent of both. The
research question is whether shared execution infrastructure can enforce
per-user and per-pool resource policies, and attribute every run to the
correct account and funding source, when the authoritative cost record (AWS
billing) arrives hours after the run. During months 13--30, the project will
develop an allocation service administered from the Control Center: funded
computing pools; introductory allocations that can be enabled, disabled, and
sized; allocations to individuals or approved groups that can be raised,
reduced, or revoked; eligible applications, workflows, and resource profiles;
job-duration and per-user concurrency limits; and an append-only record of
every allocation change and charge. Allocations are internal accounting units,
not AWS promotional credits. Before submission a run reserves its estimated
maximum cost (resource profile and duration limit plus storage and transfer
allowances); the reservation is reconciled against measured runtime and later
billing records, with explicit treatment of failed and interrupted jobs,
cleanup, and estimation error. New submissions stop when an allocation or pool
cannot cover a reservation; billing alerts remain a second line of defense,
not a limit. The current code uses the host's own AWS credentials when a user
has no connected account; the service will require an explicit execution mode
and refuse workload submission under host credentials except for authorized
administrators. Validation will test account attribution, refusal of
unauthorized and over-limit submissions, isolation between pools, recovery
after interrupted jobs, auditability of allocation changes, and estimated
versus actual cost; dollar-level limits will be claimed only to the precision
these tests support. Distributing grant-funded computing through the service
requires Georgia Tech approval.

\subsection{Reproducible experiment services and community access (\textbf{O3})}

The O3 research question is whether persistent, versioned experiment records
and application-specific evaluation make results from four different kinds of
cryosphere applications reproducible and usable by researchers outside the
development team.

\textbf{Task 3.1 -- Durable identity, experiment, and result/provenance
services.} The current password, GitHub, and ORCID identities, RBAC, saved
configurations, workspaces, experiments, events, and Control Center
demonstrate the service model in SQLite. Moving multi-user production state to
a transactional shared database with versioned migrations, encrypted backups,
retention, and export, and replacing the FastAPI relay's process-local
connector registry and pending futures with a durable task/session store, is
necessary enabling infrastructure, not itself the contribution of this task.
The contribution is a common, versioned experiment/result contract that binds,
for every qualified application and backend, the scientific/application
configuration, execution metadata, backend and environment identity,
scientific outputs, application-specific diagnostics, provenance,
software/environment versions, and citation/export information. ISSM and
Icepack already export a structured result package; ICESEE does not, and its
data-assimilation diagnostics -- ensemble error and spread, innovation
statistics, and other analysis/assimilation diagnostics appropriate to the
selected filter (for example, analysis increments or rank histograms) --
are the immediate concrete case this contract must generalize to, since a
data-assimilation experiment's scientific content lives in exactly these
diagnostics, not only in a final field. Users
will be able to clone a prior configuration, launch it on another qualified
backend, and export a citation-ready run bundle built on this contract. A
documented observation-bundle schema will connect a Frozen Legacies or LIVIST
product to CryoLauncher and ICESEE while preserving data license, uncertainty,
QC, transformations, and source provenance. The Control Center will gain
resource registration, policy/audit search, live service indicators, and
privacy-preserving metric export.

\textbf{Task 3.2 -- Usability and acceptance.} The team will conduct two
formative evaluations (end of Years 1 and 2) and one summative evaluation in
Year 3. Tasks are matched to each application:
configure and run a supported model through CryoLauncher and retrieve its
results; execute a supported ICESEE assimilation example and interpret its
error and spread diagnostics; locate a historical radar dataset in Frozen
Legacies and interpret its provenance and review status; and explore a LIVIST
englacial-temperature product and identify its inputs and scientific context.
Shared tasks are finding a tutorial and reproducing a shared run. Measures
include completion rate, time to first result, correctness of a short
interpretation item for each scientific task, critical incidents, System
Usability Scale (SUS), and qualitative barriers. Studies with human
participants will begin only after a Georgia Tech IRB determination.
Targets are at least 80\% completion of the core workflow and median SUS at
least 75 by the Year-3 evaluation. Accessibility review will cover keyboard
navigation, focus order, contrast, labels, error messages, captions, and
screen-reader checks, aiming for WCAG 2.2 AA conformance on core workflows.

\textbf{Task 3.3 -- Community and governance.} One hands-on workshop will be
offered each year, complemented by quarterly virtual office hours in Years 2
and 3. Materials will include a laptop-scale path and an HPC path. Participant
support will prioritize students, early-career researchers, minority-serving
institutions, EPSCoR jurisdictions, and institutions without local HPC staff.
A public contribution guide, roadmap, code of conduct, support policy,
security-reporting channel, adapter proposal template, and lightweight
governance document will be released in Year 1. External contributors will
receive review pairing and attribution; accepted adapters must name a
maintainer and satisfy CryoBench.

\textbf{Computing-access pilot.} As part of training and evaluation, the
project proposes a bounded pilot of project-funded AWS computing for eligible
CryoLauncher and ICESEE workflows; Frozen Legacies and LIVIST require no
computing credits. Up to 50 approved participants would receive a small
introductory allocation through the Task 2.4 service (planning value \$25
each, \$1,250 or about \$1,560 with 25\% contingency, subject to institutional
approval); larger allocations for approved scientific workflows or training
events will be estimated before they are offered. Illustrative estimates from current AWS prices put a small
default run at cents, but no measured AWS benchmark yet supports precise
per-run costs. The pilot evaluates whether controlled introductory allocations reduce
onboarding barriers, measured by first-run and workflow completion and cost
predictability, and is separate from CryoBench.
Afterwards participants continue on their own AWS accounts, PACE, or other
institutional or allocated HPC resources; external credit programs such as
CloudBank will be explored but are not assumed.

\subsection{Engineering, Release, and Acceptance Process}

Work will proceed in public issues and milestone boards. Two-week development
cycles will end in reviewed pull requests; protected branches will require
tests and review. Releases will be semantically versioned, signed where
supported, archived with DOI records, and accompanied by migration notes,
compatibility matrices, security advisories, and known scientific
limitations. A twice-yearly release train balances stability with iteration.
Acceptance decisions will combine CI results, science-lead review, security
checks, and user-evaluation evidence. Urgent security fixes may use an
expedited process with subsequent public documentation.

\subsection{Building on Existing, Recognized Capabilities}

CryoStack will reuse rather than duplicate mature investments: PISM, ISSM,
Icepack, Firedrake, PETSc, MPI, Slurm, Jupyter Book, Voilà, containers, Spack,
Georgia Tech PACE, and the existing Georgia Tech virtual-machine deployment.
The team will consult Georgia Tech research-computing, security, and
research-software specialists and the wider Science Gateways community for
usability, security, and sustainability review. CryoStack will link to upstream model documentation and
issue trackers and will contribute fixes upstream when appropriate.
The adapter SDK will also document how additional models and tools can enter
the qualification process. PISM remains the sole new production model, and one
additional Frozen Legacies collection is the sole new production dataset
integration promised by this Elements project.

\subsection{Close Collaboration and Management Plan}

\textbf{PI Brian Kyanjo} will direct the project and lead adapter design,
ICESEE integration, CryoBench, execution backends, releases, and technical
documentation. \textbf{Co-PI Alexander Robel} will lead science requirements,
PISM use-case definition, verification criteria, user evaluation, and
community engagement. Both will approve production releases. Georgia Tech
PACE/security and web-infrastructure personnel will be consulted for
deployment and incident response. PISM developers/users, historical-radar
data stewards, Science Gateway specialists, and early adopters will be invited to bounded design and
acceptance reviews; letters of collaboration will document commitments before
submission.

The team will meet weekly, review risks monthly, and conduct quarterly
milestone reviews. The PI owns scope, schedule, and technical decisions; the
Co-PI has stop authority for scientifically invalid acceptance criteria.
Unresolved decisions will be recorded in an architectural decision record and
resolved by consensus or, after one week, by the PI with rationale. The public
roadmap distinguishes committed award deliverables from aspirational features.
The major risks and mitigations are:

\begin{itemize}[leftmargin=*]
  \item \textbf{PISM integration complexity:} begin with an idealized/regional
  restartable case, freeze a minimal state/parameter scope, and validate before
  adding observations or scale.
  \item \textbf{Credential or command-injection risk:} typed allow-listed
  operations, user-held credentials, per-run isolation, quotas, logging,
  scanning, and external review precede public remote execution.
  \item \textbf{Backend variability and cloud cost:} qualify only PACE and AWS,
  retain PACE as the reference, and enforce quotas, lifecycle cleanup, and
  measured cost profiles in CryoStack itself, since AWS budget alerts notify
  but do not cap spending; pilot allocations are small and approved per
  workflow.
  \item \textbf{Historical-data heterogeneity:} use a bounded second collection,
  require manifest validation and provenance, and retain expert-reviewed QC
  rather than promising fully automatic interpretation of every radar format.
  \item \textbf{Low adoption:} recruit design partners in Year 1, measure task
  success before scaling outreach, and use PISM's existing support community
  rather than creating a competing channel.
  \item \textbf{Single-maintainer and delivery risk:} require reviews,
  runbooks, release automation, maintainer rotation, and at least two people
  able to deploy and roll back each production component. Work is ordered by
  dependency: the security baseline (Task 2.1) and CryoBench's fast tier
  precede PISM production work, and the deferrable items -- a noninteractive
  radar-processing path, adapter v2, and pilot expansion -- are named in advance.
\end{itemize}

\subsection{Integrated measurable outcomes}

\subsubsection{Deliverables}

\begin{enumerate}[leftmargin=*]
  \item Application/model/data extension kits, templates, documentation, and
  compatibility policy; production PISM adapter; second Frozen Legacies
  dataset adapter; completed modular migration of shared frontend/backend code.
  \item CryoBench validation service with public tests, reference artifacts,
  release gates, cross-backend scientific-equivalence and performance records,
  and scientific limitations.
  \item Auth-bound, allow-listed connector/relay and qualified PACE and AWS
  profiles; a policy-controlled computing-allocation service with
  reservation, reconciliation, and audit records; reproducible container/Spack
  releases with SBOMs and license manifests.
  \item Versioned experiment and provenance registry; exportable run bundles;
  observation-bundle schema and one connected Frozen Legacies- or
  LIVIST-to-model demonstration; expanded Control Center and operations export.
  \item Accessible documentation, three workshops, office hours, user studies
  covering all four applications, a computing-access pilot report (if the pilot
  is approved), contribution/governance materials, and annual public metric
  reports.
\end{enumerate}

Detailed yearly targets and collection mechanisms appear in the required
\textit{Delivery Mechanism and Community Usage Metrics} supplement. Headline
targets by the end of Year 3 are: five qualified model/application adapters
including PISM, two qualified dataset adapters, two qualified execution
profiles; at least 100 annual unique
authenticated users; at least 1,000 completed scientific runs in Year 3; eight
external adopting groups cumulatively; at least 50 annual workshop
participants; median SUS at least 75; at least three external code or
documentation contributors cumulatively; and, for the PISM production
integration, at least 90\% of changes confined to the adapter and its
declarative/configuration layer, with any remaining shared-code change
entering only through a documented extension interface. Metrics are targets
for evaluation, not claims of current adoption.

\subsubsection{Sustained and Sustainable Impacts}

The sustainability strategy has four layers. \textbf{Technical:} automated
tests, immutable environments, runbooks, backups, and low-friction releases
reduce maintenance cost. \textbf{Community:} adapters have named maintainers,
external contributions receive review support, and governance makes decisions
transparent. \textbf{Institutional:} Georgia Tech infrastructure provides the
reference gateway and HPC profile while portable recipes and user-managed AWS
prevent platform lock-in. \textbf{Scholarly:} releases, run
bundles, and documentation receive persistent identifiers and clear citation
guidance. During Year 3 the team will publish a five-year operations estimate,
identify minimum service and staffing levels, and select among institutional
support, future research awards, community allocations, and a later CSSI
Transition to Sustainability proposal based on demonstrated adoption. Core
software will remain usable locally even if the hosted gateway is unavailable.
Software sustainability is distinct from the continuing cost of computing: the
award does not fund open-ended cloud use. After the access pilot, computation
continues on user-managed AWS accounts or institutional and allocated HPC,
while the software, documentation, and data applications remain usable without
project-funded credits.

\section{Broader impacts}

CryoStack lowers barriers to advanced cryosphere computing by allowing a user
to learn on a laptop and then execute the same declared workflow on an
institutional cluster or supported user-managed cloud
deployment. Open tutorials will support graduate courses, workshops,
and self-directed learning in DA, uncertainty, reproducibility, and HPC.
Participant recruitment and support will deliberately reach students,
early-career researchers, minority-serving institutions, EPSCoR jurisdictions,
and institutions without research-software staff. Materials will use accessible
web practices and captions/alt text, and evaluation will be disaggregated only
when privacy and sample size permit responsible interpretation.

The project addresses distinct barriers. Computational access is addressed by
execution on institutional HPC and the researcher's own cloud account,
supplemented by the bounded access pilot; software complexity by packaged
environments and model adapters; data accessibility by Frozen Legacies and
LIVIST, which make historical radar observations and englacial-temperature
products discoverable with their provenance; reproducibility by versioned
experiment records and CryoBench; differences in expertise by tutorials that
span modeling, data assimilation, and data exploration; and limited
institutional infrastructure by the hosted gateway and laptop-scale tutorials.

The project will also develop the CI workforce. Contributors will learn typed
interfaces, numerical testing, secure remote execution, release engineering,
software citation, user research, and cross-disciplinary collaboration.
Permissive open-source releases and reusable adapter/validation patterns can
benefit environmental modeling communities beyond ice sheets. Ultimately,
better access to observation-constrained ice-sheet modeling strengthens the
research base used to understand sea-level change, a societal hazard affecting
coastal communities and infrastructure.

\subsection{Alignment with NSF Directorates and Divisions}

The primary alignment is \textbf{CISE/OAC}: the project translates working
software into a secure, scalable, measurable, and reusable CI service spanning
scientific data, software, and heterogeneous computing. The domain alignment is
\textbf{GEO/Office of Polar Programs}, because the scientific users and
outcomes concern ice-sheet dynamics, observations, and sea-level-relevant
cryosphere research. A secondary alignment is \textbf{GEO/EAR} through
geophysical inverse problems and Earth-system modeling. The project also
responds to CISE interest in software engineering that advances another
scientific discipline. Before submission, the team will discuss scope,
directorate fit, and the upper-range Elements budget with the cognizant OAC and
OPP program officers.

\section{Timeline}
\label{sec:timeline}

\begin{center}
\scriptsize
\begin{tabular}{@{}p{0.06\textwidth}p{0.28\textwidth}p{0.29\textwidth}p{0.27\textwidth}@{}}
\toprule
Period & O1: applications & O2: execution & O3: community platform \\
\midrule
Y1 & Adapter v1; PISM proof of concept; fast CryoBench tier & Threat model;
typed connector; PACE profile; provenance schema; license reconciliation &
Design partners; governance kit; workshop; IRB determination; formative study; baseline metrics \\
Y2 & PISM production release; second Frozen collection; observational demonstration; all adapters gated
& AWS qualification; allocation service; SBOM/scanning; run export; security review &
Workshop; office hours; second study; access pilot; external contribution \\
Y3 & Cross-model and data-to-model experiments; limitations/performance catalog; adapter v2 if
needed & Two backend profiles; allocation validation; recovery exercise; cost/operations plan;
stable release & Workshop; summative study; maintainer handoff; five-year
sustainability and final metrics \\
\bottomrule
\end{tabular}

\vspace{3pt}
\parbox{0.92\textwidth}{\footnotesize\textbf{Table 1.} Integrated work plan.
Production promotion requires scientific, engineering, security, and
user-acceptance evidence rather than schedule alone.}
\end{center}

\section{Personnel}

\textbf{PI Brian Kyanjo} is a research software developer specializing in
ensemble data assimilation, model coupling, MPI/HPC workflows, reproducible
environments, and browser-mediated scientific computing. He is the lead
developer of ICESEE and the CryoStack/CryoLauncher prototypes. He will direct
the project and lead O1, O2, release engineering, and platform integration.

\textbf{Co-PI Alexander Robel} is a Georgia Tech cryosphere scientist whose
work combines ice-sheet modeling, theory, uncertainty, and data assimilation.
He will lead scientific requirements and acceptance, the PISM use case, user
evaluation, and community engagement across O1 and O3. Georgia Tech PACE,
security, and web-infrastructure personnel, experienced PISM users,
historical-radar data stewards, research-software specialists, and early adopters will participate
through bounded design, security, and acceptance reviews documented before
submission.

\section{Results from prior NSF support}

Co-PI Alexander Robel is PI of NSF Award 2235920,
\textit{CAREER: What's Past is Prologue: Seamless Assimilation of Past
Observations into Simulations of Future Ice Sheets}, \$780,000, January 1,
2024--December 31, 2028 \citep{nsf2235920}.

\textit{Intellectual Merit:} The award develops methods and open-source
software for assimilating modern and geological observations into transient
ice-sheet simulations. ICESEE implements four ensemble-filter variants,
distributed-memory ensemble execution, adaptive state indexing, parallel
HDF5/NetCDF I/O, and state and parameter estimation. ISSM and Icepack
applications demonstrate interoperability and strong/weak scaling
\citep{kyanjo2026icesee}. That award establishes scientific DA methods; this
Elements project is distinct in creating, securing, qualifying, delivering,
and evaluating reusable infrastructure for multiple applications and
computing environments.

\textit{Broader Impacts and publication:} ICESEE is openly released, has a
citable archival version \citep{kyanjo2026zenodo}, and is distributed through
source, container, and Spack pathways. The award has supported training in
scientific software and reproducible ice-sheet modeling; its resulting
publication is \cite{kyanjo2026icesee}.
